% Generated draft: requires mathematical and source review.
\documentclass[11pt]{article}
\usepackage[margin=25mm]{geometry}
\usepackage{fontspec}
\setmainfont{Latin Modern Roman}
\usepackage{amsmath,amssymb,mathtools,unicode-math}
\setmathfont{Latin Modern Math}
\usepackage{enumitem,needspace,xurl,hyperref}
\setlength{\parindent}{0pt}
\setlength{\parskip}{5pt}
\setlength{\emergencystretch}{3em}
\newcommand{\N}{\mathbb{N}}
\newcommand{\Z}{\mathbb{Z}}
\newcommand{\Q}{\mathbb{Q}}
\newcommand{\R}{\mathbb{R}}
\newcommand{\C}{\mathbb{C}}
\newcommand{\F}{\mathbb{F}}
\newcommand{\A}{\mathbb{A}}
\newcommand{\PP}{\mathbb{P}}
\newcommand{\E}{\mathbb{E}}
\newcommand{\eps}{\varepsilon}
\newcommand{\dd}{\,\mathrm{d}}
\begin{document}
\section*{UnsolvedMath ID: 30003632}
\textbf{Source problem code:} catalogue entry
\subsection{Definitions and mathematical statement}
A group \(G\) is a \emph{Kähler group} if there exist a connected compact Kähler manifold \(X\) and a base point \(x\in X\) such that \(G\cong\pi_1(X,x)\). For this draft, use the standard smooth complex conventions: \(G\) is a \emph{projective group} if \(G\cong\pi_1(Y,y)\) for some connected smooth complex projective variety \(Y\) and \(y\in Y\); it is a \emph{quasiprojective group} if \(G\cong\pi_1(U,u)\) for some connected smooth complex quasiprojective variety \(U\) and \(u\in U\). Here quasiprojective means isomorphic to a Zariski-open subvariety of a complex projective variety, and fundamental groups are taken in the analytic topology. The source does not explicitly specify smoothness or the ground field for its algebraic realizations; these conventions therefore require confirmation against the cited report, particularly whether singular varieties are intended to be allowed.

\textbf{Classification question.} Given a specified class \(\mathcal C\) of groups, characterize each of the following subclasses up to group isomorphism:
\[
\begin{aligned}
\mathcal C_{\mathrm K}&=\{G\in\mathcal C: G\text{ is a Kähler group}\},\\
\mathcal C_{\mathrm P}&=\{G\in\mathcal C: G\text{ is a projective group}\},\\
\mathcal C_{\mathrm{QP}}&=\{G\in\mathcal C: G\text{ is a quasiprojective group}\}.
\end{aligned}
\]
Thus, for each of the three realization types, seek necessary and sufficient conditions on a member \(G\in\mathcal C\) for the corresponding realization to exist. The supplied statement requests a natural class of groups but specifies neither that class nor a meaning of naturality. This is consequently a classification-question template, rather than a fully specified classification problem for a particular class.
\subsection{Short English statement}
For a specified class of groups, determine exactly which members occur as fundamental groups of compact Kähler manifolds, smooth complex projective varieties, or smooth complex quasiprojective varieties. The source leaves the class unspecified, and the smooth complex conventions require confirmation. Its open status is supplied dated curation, not an independently verified conclusion.
\subsection{Sources}
\begin{itemize}
\item Local UnsolvedMath registry, permanent ID \(30003632\), source problem code \texttt{catalogue entry}, title: \emph{Kähler, Projective, and Quasiprojective Group Realizability}. The record belongs to the Oberwolfach Reports Open Problems collection and gives proposed year 2017; no individual proposer is supplied.
\item Supplied report citation: \emph{Mini-Workshop: Interactions between Low-dimensional Topology and Complex Algebraic Geometry} (2018), Oberwolfach Reports, pp.~2987--3033, DOI \texttt{10.4171/owr/2017/50}, \url{https://doi.org/10.4171/owr/2017/50}. The literature-source list labels the same reference \emph{Oberwolfach Report 50/2017}. The difference between the supplied citation year 2018 and the report designation 2017 is unresolved here.
\item Literature assessment checked at \texttt{2026-08-21}: supplied dated curation, not independently verified. The record labels the problem open and states that cited examples settle several individual classes of groups, while the extracted broad classification question has no general solution. Its listed checked source is \emph{Oberwolfach Report 50/2017}, \url{https://doi.org/10.4171/owr/2017/50}. The background explicitly requests independent verification of time-sensitive status and cited claims.
\end{itemize}
\end{document}
